\documentclass[11pt,a4paper]{article}

\usepackage{amsmath,amssymb,amsthm}
\usepackage{mathtools}
\usepackage{geometry}
\usepackage{hyperref}
\usepackage{booktabs}
\usepackage{array}
\usepackage{xcolor}
\usepackage{enumitem}
\usepackage{microtype}
\usepackage{float}
\usepackage{caption}
\usepackage{multirow}

\geometry{margin=2.5cm, top=2.5cm, bottom=2.5cm}

\definecolor{deepblue}{RGB}{0,51,102}
\definecolor{deepred}{RGB}{140,0,0}

\hypersetup{colorlinks=true, linkcolor=deepblue, citecolor=deepred, urlcolor=deepblue}

\newtheorem{theorem}{Theorem}[section]
\newtheorem{proposition}[theorem]{Proposition}
\theoremstyle{definition}
\newtheorem{definition}[theorem]{Definition}
\newtheorem*{remark}{Remark}

\newcommand{\dd}{\mathrm{d}}

\title{%
\textbf{Criticality in Cognitive Systems}\\[0.5em]
\large Analytic Combinatorics of Neural Avalanches,\\
Associative Memory, and Network Cascades
}

\author{
Saoirse Amarteifio\\
\small \textit{Percolation Labs}
}

\date{March 2026}

\begin{document}
\maketitle

\begin{abstract}
We apply the analytic combinatorics framework --- Lagrange equations,
Newton polygon singularity analysis, and the Flajolet--Odlyzko transfer
theorem --- to three problems at the intersection of neuroscience,
network science, and statistical physics.
(1)~\textbf{Neural avalanches}: cortical activity cascades follow power
laws $P(s) \sim s^{-\tau}$ with $\tau \approx 3/2$, consistent with the
critical brain hypothesis. We derive this from the Lagrange equation for
branching processes on cortical networks, showing that the exponent depends
on the cortical connectivity distribution via the Puiseux exponent of the
network kernel.
(2)~\textbf{Associative memory}: the Hopfield network storage capacity
transition is an absorbing-state phase transition whose order parameter
(retrieval overlap) is governed by a self-consistent equation that is a
Lagrange equation. The singularity type determines whether the transition
is continuous (second-order, square root) or discontinuous (first-order, pole).
(3)~\textbf{Cascade universality}: epidemic spreading, financial contagion,
social influence cascades, and neural avalanches are all governed by
Lagrange equations with network-dependent kernels. The AC framework
classifies them by their Puiseux exponent, revealing that the power-law
exponent $\tau$ varies continuously with the network degree distribution
exponent $\gamma$:
$\tau = (2\gamma - 3)/(\gamma - 2)$ for $3 < \gamma < 4$,
recovering mean-field $\tau = 5/2$ for $\gamma > 4$.
\end{abstract}

\tableofcontents
\clearpage

%% ================================================================
\section{Introduction}
\label{sec:intro}

The hypothesis that the brain operates near a critical point --- the
``critical brain hypothesis'' --- has been one of the most productive
ideas in computational neuroscience over the past two
decades~\cite{beggs2003,chialvo2010,munoz2018}. The experimental
evidence is striking: spontaneous cortical activity organises into
neuronal avalanches whose sizes and durations follow power
laws~\cite{beggs2003}, with exponents consistent with a branching
process at criticality.

The theoretical framework for understanding these power laws has
largely relied on two approaches: (i)~mapping to known universality
classes (directed percolation, mean-field branching
processes)~\cite{munoz2018}, and (ii)~numerical simulation of specific
neural network models~\cite{levina2007}. Both have limitations: the
universality class mapping is often argued by analogy rather than
derived from the network structure, and simulations are expensive and
model-specific.

In this paper, we introduce a third approach: \textbf{analytic
combinatorics}. The generating function for cascade sizes on any
network satisfies a Lagrange equation $T = z\,\phi(T)$, where the
kernel $\phi$ is determined by the network's degree distribution.
The singularity type of this equation --- accessible via the Newton
polygon and Puiseux expansion --- directly determines the cascade
size exponent $\tau$. No simulation is needed; no universality class
is assumed. The exponent is \emph{computed} from the network
structure.

This paper applies this framework to three systems:

\begin{enumerate}[leftmargin=2em]
\item \textbf{Neural avalanches} on cortical networks with log-normal
or truncated power-law degree distributions.

\item \textbf{Hopfield associative memory}, where the pattern retrieval
transition is reformulated as a Lagrange equation whose singularity
type determines the nature of the phase transition.

\item \textbf{Cascade universality} across domains --- epidemics, finance,
social influence, and neural cascades --- unified by the Puiseux
exponent of the network Lagrange equation.
\end{enumerate}


%% ================================================================
\section{The Lagrange Equation for Network Cascades}
\label{sec:lagrange}

\subsection{From network structure to generating function}

Consider a network with $N$ nodes and degree distribution $P(k)$
(fraction of nodes with degree $k$). Define the degree generating
function
\begin{equation}
G_0(x) = \sum_{k=0}^{\infty} P(k)\, x^k,
\end{equation}
and the \emph{excess degree} generating function
\begin{equation}
G_1(x) = \frac{G_0'(x)}{G_0'(1)} = \frac{\sum_k k\,P(k)\,x^{k-1}}{\langle k \rangle}.
\end{equation}

A cascade (avalanche, epidemic, default) is triggered by activating
a single node. Each active node activates each neighbour independently
with probability $p$ (the ``transmissibility''). The generating function
$H_1(x)$ for the size of the cascade reached via a randomly chosen edge
satisfies the self-consistent equation~\cite{newman2001}
\begin{equation}
H_1(x) = x\,\bigl(1 - p + p\,G_1(H_1(x))\bigr).
\label{eq:cascade_lagrange}
\end{equation}
This is a \textbf{Lagrange equation} $T = z\,\phi(T)$ with kernel
$\phi(T) = 1 - p + p\,G_1(T)$.

The cascade size distribution for a randomly chosen starting node is
then
\begin{equation}
H_0(x) = x\,\bigl(1 - p + p\,G_0(H_1(x))\bigr),
\end{equation}
and the probability of a cascade of size $s$ is $P(s) = [x^s]\,H_0(x)$.

\subsection{Singularity determines the exponent}

The transfer theorem~\cite{flajolet2009} states: if $H_0(x)$ has an
algebraic singularity of the form $(1 - x/x_*)^{\alpha}$ at its
dominant singularity $x_*$, then
\begin{equation}
P(s) = [x^s]\,H_0 \sim \frac{s^{-\alpha - 1}}{\Gamma(-\alpha)}\, x_*^{-s}.
\end{equation}
At the critical point (percolation threshold), $x_* = 1$ (the
singularity sits on the unit circle), and the exponential factor
disappears, leaving a pure power law:
\begin{equation}
P(s) \sim s^{-\tau}, \qquad \tau = \alpha + 1.
\end{equation}

The Puiseux exponent $\alpha$ is determined by the Newton polygon of
the implicit function $F(T, x) = T - x\,\phi(T)$ shifted to the
branch point.


%% ================================================================
\section{Neural Avalanches}
\label{sec:neural}

\subsection{The critical brain hypothesis}

Beggs and Plenz~\cite{beggs2003} discovered that spontaneous activity
in cortical slices organises into neuronal avalanches with a size
distribution $P(s) \sim s^{-3/2}$. This exponent matches the
mean-field prediction for a critical branching process, leading to the
hypothesis that the cortex operates near the critical point of a
branching process (branching ratio $\sigma \approx 1$).

\subsection{From branching ratio to Lagrange equation}

A neuron fires and activates each of its $k$ postsynaptic targets
independently with probability $p$. The total number of activations
from one firing is $\text{Binomial}(k, p)$, and the number of
\emph{secondary} activations follows the offspring generating function
\begin{equation}
G_1^{\rm eff}(x) = G_1\bigl(1 - p + p\,x\bigr).
\end{equation}
The mean offspring is $\sigma = p\,\langle k \rangle_{\rm excess}$,
the \emph{branching ratio}. At criticality, $\sigma = 1$.

The avalanche size GF satisfies
\begin{equation}
H(x) = x\,G_1^{\rm eff}\bigl(H(x)\bigr) = x\,G_1\bigl(1 - p + p\,H(x)\bigr),
\end{equation}
which is a Lagrange equation with kernel $\phi(T) = G_1(1 - p + p\,T)$.

\subsection{Exponent depends on cortical connectivity}

\paragraph{Homogeneous (Poisson) connectivity.}
For $G_1(x) = e^{c(x-1)}$, the kernel is $\phi(T) = e^{cp(T-1)}$,
which is entire. The branch point is a standard square root:
\begin{equation}
\alpha = 1/2, \qquad \tau = 3/2.
\end{equation}
This is the Beggs--Plenz exponent, consistent with the experimental
observation. The mean-field branching process corresponds to a
Poisson (or Erd\H{o}s--R\'enyi) connectivity.

\paragraph{Scale-free connectivity.}
Cortical networks may have broad degree distributions. If
$P(k) \sim k^{-\gamma}$ with $3 < \gamma < 4$, then $G_1$ has a
fractional singularity at $x = 1$ of type $(1-x)^{\gamma - 3}$.
The Lagrange equation inherits this:
\begin{equation}
\alpha = \frac{1}{\gamma - 2}, \qquad
\tau = \frac{2\gamma - 3}{\gamma - 2}.
\end{equation}
For $\gamma = 3.5$: $\tau = 8/3 \approx 2.67$, significantly larger
than $3/2$. This predicts that cortical regions with broader
connectivity distributions should have \emph{steeper} avalanche size
distributions --- a testable prediction.

\paragraph{Log-normal connectivity.}
Empirical cortical connectivity is often well fit by a log-normal
distribution~\cite{buzsaki2014}. For $P(k) \sim
\frac{1}{k}\exp\bigl(-\frac{(\ln k - \mu)^2}{2\sigma^2}\bigr)$,
$G_0(x)$ is entire (no power-law singularity). Therefore $G_1$ is
also analytic at $x = 1$, and the standard square-root branch applies:
\begin{equation}
\text{Log-normal connectivity:} \qquad \tau = 3/2.
\end{equation}
This explains why the experimental exponent $\tau \approx 3/2$ is
robust: it requires only that the degree distribution have finite
moments (no power-law tail), which is the case for log-normal and
all compact distributions.

\subsection{Prediction: when $\tau \neq 3/2$}

The AC framework makes a sharp prediction: $\tau \neq 3/2$ if and
only if the degree distribution has a power-law tail with $3 < \gamma < 4$.
This is testable: measure the degree distribution of a cortical
network (e.g.\ from connectomics data) and compute $G_1$. If $G_1$ is
analytic at $x = 1$, predict $\tau = 3/2$. If $G_1$ has a fractional
singularity, predict $\tau = (2\gamma - 3)/(\gamma - 2)$.


%% ================================================================
\section{Hopfield Networks and Associative Memory}
\label{sec:hopfield}

\subsection{The model}

The Hopfield network~\cite{hopfield1982} stores $P$ binary patterns
$\{\xi^{\mu}_i\}$ ($\mu = 1, \ldots, P$; $i = 1, \ldots, N$;
$\xi^{\mu}_i = \pm 1$) via the Hebbian coupling
\begin{equation}
J_{ij} = \frac{1}{N} \sum_{\mu=1}^{P} \xi^{\mu}_i \xi^{\mu}_j.
\end{equation}

The retrieval quality is measured by the overlap $m = \frac{1}{N}
\sum_i \xi^{1}_i \langle S_i \rangle$ with the target pattern $\xi^1$.
In the thermodynamic limit $N \to \infty$ with $\alpha = P/N$ fixed,
the overlap satisfies the self-consistent equation~\cite{amit1985}
\begin{equation}
m = \int_{-\infty}^{\infty} \frac{e^{-z^2/2}}{\sqrt{2\pi}}\,
\tanh\bigl(\beta\,m + \beta\,\sqrt{\alpha r}\,z\bigr)\,\dd z,
\label{eq:hopfield_sc}
\end{equation}
where $r$ is the spin-glass order parameter and $\beta = 1/T$ is the
inverse temperature.

\subsection{Reformulation as a Lagrange equation}

At zero temperature ($\beta \to \infty$), the self-consistent equation
simplifies to $m = \text{erf}(m / \sqrt{2\alpha})$. Define
$T = m$ and $z = 1$ (the equation is implicit in $m$ alone with
$\alpha$ as parameter):
\begin{equation}
m = \phi_\alpha(m) \equiv \text{erf}\!\left(\frac{m}{\sqrt{2\alpha}}\right).
\end{equation}

This is a fixed-point equation $T = \phi(T)$, which is the Lagrange
equation $T = z\,\phi(T)$ evaluated at $z = 1$. The transition from
retrieval ($m > 0$) to confusion ($m = 0$) occurs where $\phi'(0) = 1$,
i.e.\ $\sqrt{2/(\pi\alpha)} = 1$, giving the critical capacity
\begin{equation}
\alpha_c = 2/\pi \approx 0.637 \qquad \text{(zero temperature)}.
\end{equation}

The nature of the transition is determined by the singularity of
$m(\alpha)$ at $\alpha_c$. Expanding $\phi_\alpha(m)$ around $m = 0$:
\begin{equation}
\phi_\alpha(m) = \sqrt{\frac{2}{\pi\alpha}}\,m
- \frac{1}{3}\left(\frac{2}{\pi\alpha}\right)^{3/2} m^3 + O(m^5).
\end{equation}

At $\alpha = \alpha_c$: the linear term has coefficient $1$ (marginal),
and the cubic term has coefficient $-\frac{1}{3}(2/\pi\alpha_c)^{3/2}
= -\frac{1}{3\sqrt{\alpha_c}} < 0$. The fixed point $m^* > 0$ satisfies
\begin{equation}
m^* \approx \sqrt{3\sqrt{\alpha_c}\,(\alpha_c - \alpha)} \sim
(\alpha_c - \alpha)^{1/2}.
\end{equation}

This is a \textbf{square-root branch}: the order parameter vanishes
continuously as $m \sim (\alpha_c - \alpha)^{1/2}$, a second-order
(continuous) phase transition. The Puiseux exponent is $1/2$, the same
as the mean-field universality class.

\subsection{The full Hopfield transition at finite temperature}

At finite temperature ($\beta < \infty$), the self-consistent
equation~\eqref{eq:hopfield_sc} has a richer structure. Below the
spin-glass transition temperature $T_g$, the retrieval transition
becomes \emph{first-order} (discontinuous): $m$ jumps from a positive
value to zero as $\alpha$ increases through $\alpha_c(T)$. In the
Lagrange framework, this corresponds to the branch point moving
\emph{off the real axis} (into the complex plane), so the singularity
is no longer accessible along the real axis. Instead, the system
undergoes a \emph{spinodal} --- the metastable retrieval state
disappears.

The AC classification:
\begin{itemize}[leftmargin=2em]
\item $T = 0$: continuous transition, square-root branch, $m \sim (\alpha_c - \alpha)^{1/2}$.
\item $0 < T < T_g$: first-order transition. The Lagrange equation has
  real branch points, but the physical minimum is separated from the
  $m = 0$ state by a barrier. The system jumps discontinuously.
\item $T = T_g$: the branch point coincides with the spin-glass
  transition. Higher-order singularity (tricritical point).
\end{itemize}


%% ================================================================
\section{Cascade Universality Across Domains}
\label{sec:universality}

The Lagrange equation~\eqref{eq:cascade_lagrange} governs cascades in
many domains. The AC framework reveals that they are all classified by
the same principle: the Puiseux exponent of the network kernel $G_1$.

\begin{table}[H]
\centering
\caption{Cascade universality: the same Lagrange equation,
different kernels, different exponents.}
\renewcommand{\arraystretch}{1.2}
\begin{tabular}{p{3cm}p{3cm}ccc}
\toprule
Domain & Cascade & $G_1$ type & Puiseux & $\tau$ \\
\midrule
Cortical avalanches & Neuron firing & log-normal (analytic) & $1/2$ & $3/2$ \\
SIR epidemic & Infection spread & depends on network & varies & varies \\
Financial contagion & Bank default & interbank (heavy-tailed) & $<1/2$ & $>5/2$? \\
Social influence & Retweet/share & power-law ($\gamma \approx 3$) & $1/(\gamma{-}2)$ & $(2\gamma{-}3)/(\gamma{-}2)$ \\
Power grid & Cascading failure & grid topology (sparse) & $1/2$ & $3/2$ \\
Protein misfolding & Prion propagation & spatial (lattice) & $1/2$ & $3/2$ (DP) \\
\bottomrule
\end{tabular}
\end{table}

The key insight: the cascade exponent $\tau$ is determined entirely
by the \emph{singularity type} of the network's degree generating
function. Systems with analytic $G_1$ (finite-moment degree distributions:
Poisson, log-normal, exponential) all give $\tau = 3/2$ regardless of
details. Only power-law degree distributions with $3 < \gamma < 4$
produce anomalous exponents.

This explains the empirical regularity that $\tau \approx 3/2$ is observed
across wildly different systems --- from cortical avalanches to earthquake
aftershocks to financial crashes. It is not a coincidence or fine-tuning:
it is the \emph{generic} singularity of any Lagrange equation with an
analytic kernel.


%% ================================================================
\section{Discussion}
\label{sec:discussion}

The AC framework brings three advantages to the study of critical
phenomena in cognitive and network systems:

\paragraph{1. Computation replaces analogy.}
Traditional approaches classify neural avalanches as ``DP-like'' or
``mean-field branching'' by analogy with known universality classes.
The AC approach \emph{computes} the exponent from the network's degree
distribution --- no analogy needed. If the degree distribution changes
(e.g.\ during development or disease), the predicted exponent changes
with it.

\paragraph{2. Non-universal exponents are predicted.}
The standard critical brain hypothesis predicts $\tau = 3/2$ universally.
The AC framework predicts $\tau = 3/2$ for analytic degree distributions
but \emph{continuously varying} $\tau$ for heavy-tailed distributions.
This is testable: cortical regions with broader connectivity (e.g.\
hub-rich prefrontal cortex vs.\ regular sensory cortex) should show
different avalanche exponents.

\paragraph{3. Phase transition type is classified.}
The Hopfield network retrieval transition is classified by its
singularity: continuous at $T = 0$ (square-root branch), first-order at
$T > 0$ (complex branch point), tricritical at $T = T_g$ (higher-order
singularity). This classification extends to any associative memory
model whose order parameter satisfies a self-consistent (Lagrange)
equation.


%% ================================================================
\section{Conclusion}
\label{sec:conclusion}

Analytic combinatorics provides a unified, computational framework for
critical phenomena in cognitive and network systems. The singularity type
of the Lagrange equation --- determined by the network's degree
distribution --- is the single organising principle that classifies
cascade exponents, distinguishes continuous from discontinuous phase
transitions, and predicts when the ``universal'' exponent $\tau = 3/2$
should and should not hold.

The framework is directly applicable to empirical data: given a measured
degree distribution (from connectomics, social network data, or financial
network data), the pipeline computes $G_0 \to G_1 \to$ Lagrange equation
$\to$ Newton polygon $\to$ Puiseux exponent $\to$ $\tau$. No simulation
is required.


%% ================================================================
\begin{thebibliography}{99}

\bibitem{beggs2003}
J.~M.~Beggs and D.~Plenz,
\newblock Neuronal avalanches in neocortical circuits,
\newblock \emph{J.\ Neurosci.}\ \textbf{23}, 11167--11177 (2003).

\bibitem{chialvo2010}
D.~R.~Chialvo,
\newblock Emergent complex neural dynamics,
\newblock \emph{Nature Physics}\ \textbf{6}, 744--750 (2010).

\bibitem{munoz2018}
M.~A.~Mu\~noz,
\newblock Colloquium: Criticality and dynamical scaling in living systems,
\newblock \emph{Rev.\ Mod.\ Phys.}\ \textbf{90}, 031001 (2018).

\bibitem{levina2007}
A.~Levina, J.~M.~Herrmann, and T.~Geisel,
\newblock Dynamical synapses causing self-organized criticality in neural
networks,
\newblock \emph{Nature Physics}\ \textbf{3}, 857--860 (2007).

\bibitem{hopfield1982}
J.~J.~Hopfield,
\newblock Neural networks and physical systems with emergent collective
computational abilities,
\newblock \emph{Proc.\ Natl.\ Acad.\ Sci.}\ \textbf{79}, 2554--2558 (1982).

\bibitem{amit1985}
D.~J.~Amit, H.~Gutfreund, and H.~Sompolinsky,
\newblock Spin-glass models of neural networks,
\newblock \emph{Phys.\ Rev.\ A}\ \textbf{32}, 1007--1018 (1985).

\bibitem{newman2001}
M.~E.~J.~Newman, S.~H.~Strogatz, and D.~J.~Watts,
\newblock Random graphs with arbitrary degree distributions and their
applications,
\newblock \emph{Phys.\ Rev.\ E}\ \textbf{64}, 026118 (2001).

\bibitem{flajolet2009}
P.~Flajolet and R.~Sedgewick,
\newblock \emph{Analytic Combinatorics}, Cambridge University Press (2009).

\bibitem{buzsaki2014}
G.~Buzs\'aki and K.~Mizuseki,
\newblock The log-dynamic brain: how skewed distributions affect network
operations,
\newblock \emph{Nature Rev.\ Neurosci.}\ \textbf{15}, 264--278 (2014).

\bibitem{cohen2002}
R.~Cohen, D.~ben-Avraham, and S.~Havlin,
\newblock Percolation critical exponents in scale-free networks,
\newblock \emph{Phys.\ Rev.\ E}\ \textbf{66}, 036113 (2002).

\end{thebibliography}

\end{document}
